%====================================================================
%  TidalLanes_Model_v2.tex
%
%  Journal-style condensed version of the model and estimation.
%  Derivations are in the appendix; the main text states
%  definitions, equations, and results.
%====================================================================
\documentclass[12pt]{article}

\usepackage[margin=1in]{geometry}
\usepackage{amsmath,amssymb,amsthm,mathtools}
\usepackage{graphicx,booktabs,tabularx,array,setspace,float,enumitem}
\usepackage{natbib}
\usepackage[colorlinks=true,linkcolor=blue,citecolor=blue,urlcolor=blue]{hyperref}
\setstretch{1.1}
\setlength{\parindent}{1.5em}

\newtheorem{proposition}{Proposition}
\newtheorem{lemma}{Lemma}
\newtheorem{definition}{Definition}
\newtheorem{assumption}{Assumption}

\DeclareMathOperator*{\argmax}{arg\,max}

\newcommand{\N}{\mathcal N}
\newcommand{\E}{\mathcal E}
\newcommand{\R}{\mathcal R}
\newcommand{\Lbar}{\bar L}
\newcommand{\Wbar}{\bar W}
\newcommand{\ubar}{\bar u}
\newcommand{\Abar}{\bar A}
\newcommand{\tbar}{\bar t}
\newcommand{\tkl}{t_{kl}}
\newcommand{\tijk}{\tau_{ij}}
\newcommand{\obs}{\mathrm{obs}}
\newcommand{\cf}{\mathrm{cf}}

\title{Directional Congestion and Tidal Lanes}
\author{Yige Duan \and Xiaoruo Feng \and Yizhen Gu}
\date{\today}

\begin{document}
\maketitle

\begin{abstract}
\noindent
We develop a quantitative spatial model of commuting in which every
link of the transport network carries \emph{directional} primitives:
lane capacity, free-flow speed, and realised travel time differ
between the two orientations of a corridor.  Directional travel times
generate an asymmetric bilateral commuting cost matrix, asymmetric
link-traffic gravity, and direction-specific congestion externalities.
The model nests the symmetric case as a special point and delivers a
closed-form welfare representation from which the gains of tidal-lane
policies---reallocations of lane capacity across the two directions
of a corridor---can be read off directly.  We bring the model to
Chinese morning-peak commuting microdata and show that lane-preserving
tidal-lane reallocations on a small set of corridors deliver welfare
gains comparable to much larger symmetric capacity expansions.
\end{abstract}

%====================================================================
\section{Introduction}\label{sec:intro}
%====================================================================

Most quantitative urban commuting models treat corridors as
undirected: a lane added to corridor $c$ reduces travel time from $k$
to $l$ by the same amount as from $l$ to $k$.  This abstraction
understates the rigidity of real-world road networks, where peak-hour
flows run sharply in one direction and off-peak capacity on the other
is underutilised.  The policy response in many cities---notably
tidal-lane (``contraflow'') reallocations, which shift one or more
lanes toward the peak direction during peak hours---operates on an
object that undirected models cannot see.

This paper develops a directional quantitative spatial framework that
makes this object visible.  Residents choose a home, a workplace, and
a route; utility is deflated by an iceberg route cost; congestion on
each direction of each link is driven by traffic on that direction
alone; and bilateral commuting costs, link traffic, welfare, and
hat-algebra counterfactuals inherit the underlying directionality.
Our framework embeds the Fr\'echet route-choice structure of
\citet{AllenArkolakis2022} but relaxes their symmetric-link
assumption throughout; it also extends the commuting-and-agglomeration
structure of \citet{Ahlfeldt2015} to admit network-level asymmetry.

Bringing the model to morning-peak commuting data for a large Chinese
metropolitan area, we estimate the strength of directional congestion
from a speed-on-traffic gravity regression, invert residential
amenities and workplace productivities from the observed
origin--destination matrix, and evaluate counterfactual lane
reallocations.  Tidal-lane reallocations on a small set of identified
corridors deliver welfare gains that are of the same order as, but
obtained at a fraction of the cost of, uniform lane expansion.

%====================================================================
\section{Model}\label{sec:model}
%====================================================================

\subsection{Environment}\label{ssec:env}

The city consists of $N$ spatial cells $\N=\{1,\dots,N\}$ arranged on
a directed graph $(\N,\E)$ with edge set
$\E\subseteq\N\times\N$.  A directed edge $(k,l)\in\E$ indicates an
admissible move from $k$ to $l$; $(k,l)\in\E$ does not imply
$(l,k)\in\E$, and even when both hold the two directions are distinct
objects.  Each directed link $(k,l)$ carries three primitives:
physical length $\ell_{kl}>0$, lane capacity $n_{kl}>0$, and free-flow
speed $v^{0}_{kl}>0$, with free-flow travel time
$\tbar_{kl}\equiv\ell_{kl}/v^{0}_{kl}$.  The realised (congested)
directed travel time $\tkl>0$ is determined in equilibrium below.

The city is populated by a unit mass $\Lbar$ of workers, allocated
across residences $\{L_i^R\}$ and workplaces $\{L_j^F\}$ with
$\sum_iL_i^R=\sum_jL_j^F=\Lbar$.  Each cell has a baseline residential
amenity $\ubar_i>0$ and workplace productivity $\Abar_j>0$.  Endogenous
amenities and productivities depend on the local scale of activity,
\begin{equation}
  u_i=\ubar_i(L_i^R)^{\beta},\qquad
  A_j=\Abar_j(L_j^F)^{\alpha},\qquad
  \alpha>0,\ \beta<0,
  \label{eq:fundamentals}
\end{equation}
where $\alpha$ is the (net) workplace agglomeration elasticity and
$\beta$ is the (net) residential congestion elasticity.  The model is
closed by four structural parameters: $\theta>0$ governing the
dispersion of route choice, $(\alpha,\beta)\in\mathbb R^2$ in
\eqref{eq:fundamentals}, and the link-level congestion elasticity
$\lambda\ge 0$ introduced below.

\subsection{Route choice and bilateral commuting cost}\label{ssec:route}

A worker $\nu$ chooses a residence $i$, a workplace $j$, and a
commuting route $r=(r_0=i,r_1,\dots,r_K=j)$ with $(r_{m-1},r_m)\in\E$.
Indirect utility is
\begin{equation}
  V_{ij,r}(\nu)
  \;=\;
  \Bigl(\frac{u_iw_j}{\prod_{m=1}^{K}t_{r_{m-1},r_{m}}}\Bigr)\varepsilon_{ij,r}(\nu),
  \qquad
  \varepsilon_{ij,r}\overset{\text{iid}}{\sim}\text{Fr\'echet}(\theta),
  \label{eq:utility}
\end{equation}
where $w_j$ is the workplace wage and $\varepsilon$ is an
idiosyncratic taste shock with unit scale.  Consistent with the
quantitative spatial literature, we treat each link cost $\tkl$ as a
\emph{dimensionless iceberg index} with $\tkl\ge 1$: route costs are
multiplicative in $\tkl$ and changes in the time unit (seconds,
minutes, hours) rescale every $\tkl$ by a common factor that is
absorbed by the aggregate normalisation introduced in
\S\ref{ssec:cong}.\footnote{The micro-foundation, developed formally in
Appendix~\ref{app:units}, is $\tkl=(\ell_{kl}/v_{kl})^{1/\theta}$
with congested directed speed
$v_{kl}=v^{0}_{kl}(n_{kl}/\xi_{kl})^{\lambda\theta}$.  The choice
$\delta_0=1/\theta$ delivers a gravity-consistent distance elasticity
of $-1$.}  The link costs $\{\tkl\}$ are endogenous and determined by
the congestion technology of \S\ref{ssec:cong}.

\paragraph{Bilateral cost index.}
Summing over feasible routes yields the \emph{bilateral commuting
cost index}
\begin{equation}
  \tijk^{-\theta}
  \;\equiv\;
  \sum_{r\in\R_{ij}}\prod_{m=1}^{K(r)}t_{r_{m-1},r_{m}}^{-\theta}.
  \label{eq:tau_def}
\end{equation}
Collect link costs in the directed weighted adjacency matrix
$A=[a_{kl}]$ with $a_{kl}=\tkl^{-\theta}\mathbf 1\{(k,l)\in\E\}$.
When the spectral radius $\rho(A)<1$ the series
$\sum_{K\ge 0}A^{K}$ converges and
\begin{equation}
  \tijk^{-\theta}\;=\;[(I-A)^{-1}]_{ij},
  \label{eq:tau_series}
\end{equation}
so $\tijk$ is determined by the full network of directed link costs
and is generically asymmetric, $\tijk\neq\tau_{ji}$.

\begin{assumption}[Spectral radius]\label{ass:specrad}
$\rho(A)<1$, equivalently
$\sum_{l}\tkl^{-\theta}\mathbf 1\{(k,l)\in\E\}<1$ for every
$k\in\N$.
\end{assumption}

\subsection{Gravity, market access, and link traffic}\label{ssec:gravity}

Standard properties of i.i.d.\ Fr\'echet draws deliver the joint
residence--workplace--route choice probability and, upon aggregation
over routes, the commuting gravity equation.  Let
\begin{equation}
  \Wbar\;\equiv\;\Bigl[\sum_{(i,j)\in\N^2}\tijk^{-\theta}u_i^{\theta}w_j^{\theta}\Bigr]^{1/\theta}
  \label{eq:Wbar_def}
\end{equation}
denote the aggregate commuter-market-access scale.  Then
\begin{equation}
  L_{ij}\;=\;\Lbar\cdot\frac{u_i^{\theta}w_j^{\theta}\tijk^{-\theta}}{\Wbar^{\theta}},
  \qquad
  L_i^R\;=\;\sum_{j}L_{ij},\qquad L_j^F\;=\;\sum_{i}L_{ij}.
  \label{eq:gravity}
\end{equation}
Define the resident and firm market-access indices
\begin{equation}
  \Pi_i^{-\theta}
  \;\equiv\;\sum_{j}\tijk^{-\theta}L_j^FP_j^{\theta},
  \qquad
  P_j^{-\theta}
  \;\equiv\;\sum_{i}\tijk^{-\theta}L_i^R\Pi_i^{\theta},
  \label{eq:MA_def}
\end{equation}
so that the commuting gravity equation admits the market-access form
\begin{equation}
  L_{ij}\;=\;\tijk^{-\theta}\cdot\frac{L_i^R}{\Pi_i^{-\theta}}\cdot\frac{L_j^F}{P_j^{-\theta}}.
  \label{eq:gravity_MA}
\end{equation}
The directional character of $\tijk^{-\theta}$ implies
$\Pi_i\ne P_i$ in general: residents in cell $i$ have different
accessibility to the workplaces they face than firms in $i$ have to
the residents who face them.

\paragraph{Link traffic.}
Directed link traffic $\xi_{kl}$ is the expected number of traversals
of edge $(k,l)$ summed across commuting pairs.  Appendix~\ref{app:linkgravity}
establishes:

\begin{proposition}[Directional link-traffic gravity]\label{prop:linkgravity}
Under Assumption~\ref{ass:specrad},
\begin{equation}
  \xi_{kl}\;=\;\tkl^{-\theta}\,P_k^{-\theta}\,\Pi_l^{-\theta}.
  \label{eq:xi_gravity}
\end{equation}
\end{proposition}

Equation \eqref{eq:xi_gravity} is the directional analogue of the
link-traffic gravity relation and is the central link between the
spatial distribution of activity and the network structure of
traffic.  Because every component is directional, $\xi_{kl}\ne\xi_{lk}$
in general.

\subsection{Congestion technology and equilibrium}\label{ssec:cong}

Realised directed travel time is driven by directed traffic per
directed lane:
\begin{equation}
  \tkl\;=\;\tbar_{kl}\Bigl(\frac{\xi_{kl}}{n_{kl}}\Bigr)^{\!\lambda},
  \qquad \lambda\ge 0.
  \label{eq:cong}
\end{equation}
Substituting \eqref{eq:xi_gravity} into \eqref{eq:cong} and solving
for $\tkl$ delivers the reduced form
\begin{equation}
  \tkl
  \;=\;
  \tbar_{kl}^{1/(1+\theta\lambda)}\cdot
  n_{kl}^{-\lambda/(1+\theta\lambda)}\cdot
  (P_k\Pi_l)^{-\theta\lambda/(1+\theta\lambda)}.
  \label{eq:tkl_reduced}
\end{equation}
Adding a directed lane to $(k,l)$ therefore reduces the directed
travel time with elasticity $-\lambda/(1+\theta\lambda)<0$ and
increases directed traffic with elasticity $\theta\lambda/
(1+\theta\lambda)>0$.  The absolute ratio of the two elasticities
equals $\theta$: traffic responds more elastically to lane additions
than does travel time, recovering a structural form of the
``fundamental law of road congestion'' of
\citet{DurantonTurner2011}.

\emph{Tidal-lane asymmetry.}  A tidal-lane policy shifts lanes from
the off-peak to the peak direction on a set of corridors
$\E^\star\subseteq\E$ while preserving total corridor capacity
$n_{kl}+n_{lk}$.  Because each direction enters \eqref{eq:cong}
separately, the policy strictly lowers $\tkl$ and strictly raises
$t_{lk}$ on treated corridors, generating asymmetric changes in
$\{\tijk\}$, $\{L_{ij}\}$, and $(L^R,L^F)$.  Within an \emph{undirected}
model this policy is vacuous.

\paragraph{Equilibrium in population shares.}
Let $l_i^R\equiv L_i^R/\Lbar$, $l_j^F\equiv L_j^F/\Lbar$.  Combining
\eqref{eq:fundamentals}, \eqref{eq:gravity}, and market clearing
(Appendix~\ref{app:eqshares}) gives the equilibrium system
\begin{align}
  (l_i^R)^{1-\beta\theta}
  &=\chi\sum_{j}\tijk^{-\theta}\ubar_i^{\theta}\Abar_j^{\theta}(l_j^F)^{\alpha\theta},
  \label{eq:eq_R}\\
  (l_j^F)^{1-\alpha\theta}
  &=\chi\sum_{i}\tijk^{-\theta}\ubar_i^{\theta}\Abar_j^{\theta}(l_i^R)^{\beta\theta},
  \label{eq:eq_F}
\end{align}
with $\chi\equiv\Lbar^{\theta(\alpha+\beta)}/\Wbar^{\theta}$ pinned
down by normalisation.

\begin{definition}[Equilibrium]\label{def:eq}
Given $\{(\ubar_i,\Abar_j)\}$, $\{(\ell_{kl},v^{0}_{kl},n_{kl})\}$,
and $(\theta,\alpha,\beta,\lambda)$, an equilibrium is a tuple
$\bigl(\{L_i^R,L_j^F\},\{L_{ij}\},\{\tijk\},\{\tkl\},\{\xi_{kl}\}\bigr)$
satisfying \eqref{eq:tau_series}, \eqref{eq:gravity_MA},
\eqref{eq:xi_gravity}, \eqref{eq:cong}, and
\eqref{eq:eq_R}--\eqref{eq:eq_F}.
\end{definition}

\begin{proposition}[Existence and uniqueness]\label{prop:exist}
Under Assumption~\ref{ass:specrad} and $\lambda\ge 0$, an equilibrium
exists.  It is unique if
\begin{equation}
  \alpha\;\le\;\tfrac{1}{2}(1/\theta-\lambda),
  \qquad
  \beta\;\le\;\tfrac{1}{2}(1/\theta-\lambda),
  \label{eq:unique}
\end{equation}
jointly.  With $\alpha>0$, $\beta<0$, and $\theta\lambda<1$ in our
baseline calibration, the left-hand inequality in \eqref{eq:unique}
is binding and the right-hand inequality is automatic.
\end{proposition}

The existence argument extends the Brouwer-based construction of
\citet{AllenArkolakis2022}; the uniqueness bound follows their
Proposition~1(b).  Proof sketches and the full numerical verification
procedure are in Appendix~\ref{app:exist}.

\subsection{Welfare}\label{ssec:welfare}

Welfare is the ex-ante expected maximum utility.  Because the
$\varepsilon_{ij,r}$ are i.i.d.\ Fr\'echet with scale $\omega_{ij,r}=
u_iw_j/\prod_m t_{r_{m-1},r_m}$, the maximum
$M=\max_{(i,j,r)}V_{ij,r}$ is itself Fr\'echet with scale
$\Wbar$ (Appendix~\ref{app:welfare}).  Hence
\begin{equation}
  \mathbb W
  \;\equiv\;\mathbb E\!\bigl[M\bigr]
  \;=\;\Gamma(1-1/\theta)\,\Wbar
  \;=\;\Gamma(1-1/\theta)\Bigl[\sum_{(i,j)}\tijk^{-\theta}u_i^{\theta}w_j^{\theta}\Bigr]^{1/\theta},
  \label{eq:welfare_main}
\end{equation}
which admits the equivalent market-access forms
\begin{equation}
  \Wbar^{\theta}
  \;=\;\sum_{i}u_i^{\theta}\Pi_i^{-\theta}
  \;=\;\sum_{j}w_j^{\theta}P_j^{-\theta}.
  \label{eq:welfare_MA}
\end{equation}
A first-order expansion of $\ln\mathbb W$ gives the exact welfare
decomposition
\begin{equation}
  \mathrm d\ln\mathbb W
  \;=\;\sum_{i}\Lambda_i^R\,\mathrm d\ln u_i
  \;+\;\sum_{j}\Lambda_j^F\,\mathrm d\ln w_j
  \;-\;\sum_{(i,j)}\Lambda_{ij}\,\mathrm d\ln\tijk,
  \label{eq:welfare_decomp}
\end{equation}
where the weights $\Lambda_i^R=L_i^R/\Lbar$, $\Lambda_j^F=L_j^F/\Lbar$,
and $\Lambda_{ij}=L_{ij}/\Lbar$ are baseline commuting shares.  The
directional content is concentrated in the access channel: peak-hour
flows are predominantly suburb-to-CBD, so the
marginal welfare value of a reduction in inbound peak travel time
generally exceeds the marginal welfare cost of an equal rise in
outbound off-peak travel time.  Quantifying this difference is the
central counterfactual exercise of Section~\ref{sec:cf}.

\paragraph{Hat algebra.}
For counterfactual work it is convenient to work in relative changes
$\hat x\equiv x'/x$.  Using \eqref{eq:welfare_main} together with the
gravity equation to eliminate $u_i^{\theta}w_j^{\theta}\tijk^{-\theta}$
in favour of observed commuting shares
$\pi_{ij}^{\obs}\equiv L_{ij}^{\obs}/\Lbar$ gives the exact-hat
welfare formula
\begin{equation}
  \hat{\mathbb W}^{\theta}
  \;=\;\hat{\Wbar}^{\theta}
  \;=\;\sum_{(i,j)}\pi_{ij}^{\obs}\,\hat u_i^{\theta}\hat w_j^{\theta}\hat\tau_{ij}^{-\theta}.
  \label{eq:welfare_hat}
\end{equation}
Under the spillovers \eqref{eq:fundamentals},
$\hat u_i=(\hat l_i^R)^{\beta}$ and
$\hat w_j=(\hat l_j^F)^{\alpha}$, so the welfare hat depends only on
observables and on the hat system of \S\ref{sec:est} below.  A
uniform $\hat\tau_{ij}=\kappa$ for all $(i,j)$ yields
$\hat{\mathbb W}=\kappa^{-1}$, so welfare changes admit a
``time-equivalent variation'' reading: a counterfactual with
$\hat{\mathbb W}=1.02$ is equivalent to a 1.96\% reduction in every
bilateral commuting time.

%====================================================================
\section{Estimation}\label{sec:est}
%====================================================================

\subsection{Data and calibration}\label{ssec:data}

\paragraph{Data.}
The empirical analysis uses a grid-level panel for a major Chinese
metropolitan area.  The origin--destination matrix $L_{ij}^{\obs}$ is
constructed from Baidu smartphone-tracked home--work pairs aggregated
to 3~km grid cells and restricted to the morning peak
(07:00--09:00).  Directed speeds $v_{kl}^{\obs}$ and free-flow
speeds $v^{0}_{kl}$ are length-weighted harmonic averages of
centerline-level speed observations during the AM peak and the
22:00--05:00 window, respectively, for directed links supported by
matched centreline data.  Directed lane counts $n_{kl}$ are obtained
from a roadtype lookup applied to the same centrelines.  Observed
directed travel time is $\tkl^{\obs}=\ell_{kl}/v_{kl}^{\obs}$, free-flow
time is $\tbar_{kl}=\ell_{kl}/v^{0}_{kl}$, and the bilateral commuting
cost index is reconstructed via the series representation
\eqref{eq:tau_series} evaluated at $A^{\obs}$ with entries
$(\tkl^{\obs})^{-\theta}$.  Details and descriptive statistics are in
Appendix~\ref{app:data}.

\paragraph{Calibrated parameters.}
We calibrate $(\theta,\alpha,\beta)$ to central values from the
quantitative spatial literature:
\begin{equation}
  \theta\;=\;6.83,\qquad
  \alpha\;=\;0.06,\qquad
  \beta\;=\;-0.30.
  \label{eq:calib}
\end{equation}
The dispersion parameter is taken from the Berlin gravity estimates
of \citet{Ahlfeldt2015}; the agglomeration elasticity is in the
$[0.03,0.08]$ range reviewed by \citet{CombesGobillon2015} and
\citet{RosenthalStrange2004}; the residential congestion elasticity
is at the magnitude reported in \citet{MonteReddingRossiHansberg2018}
and \citet{HeblichReddingSturm2020}.  The uniqueness bound
\eqref{eq:unique} requires $\alpha\le(1/\theta-\lambda)/2\approx
0.061$ at our estimated $\lambda$ (below), which holds.  Robustness
grids $\alpha\in\{0.04,0.06,0.08\}$ and
$\beta\in\{-0.50,-0.30,-0.20\}$ are reported alongside the main
estimates.

\subsection{Congestion elasticity}\label{ssec:lambda}

The congestion technology \eqref{eq:cong}, combined with the iceberg
mapping $\tkl=(\ell_{kl}/v_{kl})^{1/\theta}$ and the definition of
$\tbar_{kl}$, yields the reduced-form speed--flow relationship
\begin{equation}
  \ln v_{kl}\;=\;\ln v^{0}_{kl}
  -\delta_1\ln\!\Bigl(\frac{\xi_{kl}}{n_{kl}}\Bigr),
  \qquad
  \delta_1\;\equiv\;\theta\lambda,
  \label{eq:speed_flow}
\end{equation}
which we estimate on the cross-section of directed links after
absorbing $k$- and $l$-fixed effects.  OLS is biased downward because
the error $\ln v^{0}_{kl}$, through the gravity equation
\eqref{eq:xi_gravity}, is negatively correlated with equilibrium
traffic $\xi_{kl}$.  We instrument $\ln(\xi_{kl}/n_{kl})$ with
observables that shift the \emph{demand} for traversing $(k,l)$ but
are uncorrelated with the free-flow component of the travel-time
cost; our baseline instrument is the count of signalised intersections
along a link conditional on bilateral route quality, following the
urban-variant identification strategy in \citet{AllenArkolakis2022}.
Table~\ref{tab:lambda} reports the results.  The baseline IV estimate
is $\hat\delta_1^{\mathrm{IV}}=0.49$ (s.e.\ 0.14), implying
$\hat\lambda=\hat\delta_1/\theta\approx0.072$, so $\theta\hat\lambda
<1$ and the uniqueness condition \eqref{eq:unique} is satisfied.
Appendix~\ref{app:iv} reports first-stage diagnostics, alternative
instruments, and an identification discussion.

\begin{table}[H]
\centering
\caption{Estimating the congestion elasticity}
\label{tab:lambda}
\begin{tabular}{lcccc}
\toprule
 & (1) & (2) & (3) & (4) \\
 & OLS & OLS & IV: 1st stage & IV: 2nd stage\\
\midrule
log traffic per lane & 0.06 & 0.03 & & 0.49\\
 & (0.008) & (0.010) & & (0.14)\\
log intersections & & & $-0.34$ & \\
 & & & (0.06) & \\
\midrule
Origin FE   & No  & Yes & Yes & Yes\\
Destination FE & No & Yes & Yes & Yes\\
F-stat (first stage) & & & 32.4 & \\
Obs.\ (directed links) & 1{,}338 & 1{,}338 & 1{,}338 & 1{,}338\\
\bottomrule
\end{tabular}
\end{table}

\subsection{Inversion of fundamentals}\label{ssec:invert}

With $(\theta,\alpha,\beta,\lambda)$ in hand we invert the equilibrium
system \eqref{eq:eq_R}--\eqref{eq:eq_F} to recover the baseline
amenity and productivity profiles $(\ubar_i,\Abar_j)$ that rationalise
the observed shares $(l_i^{R,\obs},l_j^{F,\obs})$ and bilateral cost
index $(\tijk^{\obs})$, subject to the normalisations
$\prod_i\ubar_i=\prod_j\Abar_j=1$.  The inversion alternates two
closed-form updates,
\begin{equation}
  \ubar_i^{(k+1),\theta}
  \;=\;
  \frac{(l_i^{R,\obs})^{1-\beta\theta}}
       {\chi^{(k)}\sum_{j}(\tijk^{\obs})^{-\theta}\Abar_j^{(k),\theta}(l_j^{F,\obs})^{\alpha\theta}},
  \qquad
  \Abar_j^{(k+1),\theta}
  \;=\;
  \frac{(l_j^{F,\obs})^{1-\alpha\theta}}
       {\chi^{(k)}\sum_{i}(\tijk^{\obs})^{-\theta}\ubar_i^{(k+1),\theta}(l_i^{R,\obs})^{\beta\theta}}.
  \label{eq:invert}
\end{equation}
Under Proposition~\ref{prop:exist} the level system admits a unique
strictly positive solution; we verify numerically that the Jacobian
of the fixed-point map has spectral radius strictly less than one at
each iterate (sufficient for local linear convergence) and check
invariance to the initialisation.  Convergence is reached in fewer
than 50 iterations for the 3~km grid.

\subsection{Counterfactuals}\label{sec:cf}

Counterfactuals are implemented in exact-hat algebra.  Given a
counterfactual infrastructure network
$\hat{\tbar}_{kl}=\tbar_{kl}^{\cf}/\tbar_{kl}$, $\hat n_{kl}$, the
counterfactual congested link costs follow from taking hats in
\eqref{eq:tkl_reduced}:
\begin{equation}
  \hat\tkl^{\,1+\theta\lambda}
  \;=\;\hat{\tbar}_{kl}\cdot\hat n_{kl}^{-\lambda}\cdot
       (\hat P_k\hat\Pi_l)^{-\theta\lambda};
  \label{eq:hat_tkl}
\end{equation}
$\hat\tau_{ij}$ follows from a hat version of the series
representation \eqref{eq:tau_series}; and the residential and
workplace shares satisfy the hat system
\begin{equation}
  (\hat l_i^R)^{1-\beta\theta}
  \;=\;\sum_{j}\pi_{ij|i}^{\obs}\,\hat\tau_{ij}^{-\theta}(\hat l_j^F)^{\alpha\theta},
  \qquad
  (\hat l_j^F)^{1-\alpha\theta}
  \;=\;\sum_{i}\pi_{ij|j}^{\obs}\,\hat\tau_{ij}^{-\theta}(\hat l_i^R)^{\beta\theta},
  \label{eq:hat_system}
\end{equation}
where $\pi_{ij|i}^{\obs}=L_{ij}/L_i^R$ and
$\pi_{ij|j}^{\obs}=L_{ij}/L_j^F$ are baseline conditional commuting
shares.  The welfare hat $\hat{\mathbb W}$ is read off from
\eqref{eq:welfare_hat}.  Implementation details, including the nested
fixed-point solver for the coupled system
\eqref{eq:xi_gravity}--\eqref{eq:hat_tkl}--\eqref{eq:hat_system}, are
in Appendix~\ref{app:hat}.

\paragraph{Scenarios.}
We evaluate four counterfactuals:
(CF1) recomputing the equilibrium after symmetrising
$\{\tbar_{kl}\}$ across directions, which isolates the positive
contribution of observed directionality;
(CF2) tidal-lane reallocation along a small set of identified peak
corridors, holding total lane-kilometres fixed;
(CF3) uniform 10\% lane expansion across all links; and
(CF4) a head-to-head welfare comparison of (CF2) and (CF3) at a
common investment cost.  The main finding is that the targeted
tidal-lane policy (CF2) delivers welfare gains of comparable order
to, but at a fraction of the cost of, uniform expansion (CF3); the
full counterfactual table and robustness rows are deferred to the
empirical supplement.

%====================================================================
\appendix
\section*{Appendix}
\renewcommand{\thesubsection}{\Alph{subsection}}
\setcounter{subsection}{0}
%====================================================================

\subsection{Route aggregation and the series representation}\label{app:tau}

Enumerating routes by length and using
$a_{kl}=\tkl^{-\theta}\mathbf 1\{(k,l)\in\E\}$ gives
\[
  \tijk^{-\theta}
  \;=\;\sum_{K=0}^{\infty}\sum_{r:K(r)=K}\prod_{m=1}^{K}t_{r_{m-1},r_m}^{-\theta}
  \;=\;\sum_{K=0}^{\infty}[A^K]_{ij}.
\]
When $\rho(A)<1$ the matrix geometric series converges absolutely to
$(I-A)^{-1}$, giving \eqref{eq:tau_series}.  Length-zero walks
contribute the identity diagonal and correspond to not traversing any
edge, consistent with $\tau_{ii}=1$.

\subsection{Units and the iceberg calibration}\label{app:units}

We treat $\tkl\ge 1$ as a dimensionless iceberg loss factor; route
costs are multiplicative in $\tkl$ and enter utility
\eqref{eq:utility} via $\prod_m t_{r_{m-1},r_m}$.  The micro-foundation
connecting $\tkl$ to physical travel time takes the log-linear form
$\tkl=(\ell_{kl}/v_{kl})^{\delta_0}$ with
$v_{kl}=v^{0}_{kl}(n_{kl}/\xi_{kl})^{\lambda/\delta_0}$.  Setting
$\delta_0=1/\theta$ has two consequences.  First, in the commuting
kernel $\tkl^{-\theta}$ the time unit enters only via
$(\ell_{kl}/v_{kl})^{-1}$, which delivers a gravity-consistent
distance elasticity of $-1$.  Second, a uniform change of time units
$t\mapsto ct$ rescales every $\tkl$ by a common factor $c^{1/\theta}$,
which is absorbed by the aggregate scalar $\chi$ in \S\ref{ssec:cong}
and leaves $\{L_i^R,L_j^F,L_{ij},\xi_{kl}\}$ unchanged.

\subsection{Max-of-Fr\'echet and the welfare index}\label{app:welfare}

Let $V_{ij,r}=\omega_{ij,r}\varepsilon_{ij,r}$ with
$\omega_{ij,r}=u_iw_j/\prod_m t_{r_{m-1},r_m}$.  Because
$\omega_{ij,r}\varepsilon_{ij,r}$ is Fr\'echet with shape $\theta$ and
scale $\omega_{ij,r}$, independence across $(i,j,r)$ gives
\[
  \Pr(M\le y)
  \;=\;\prod_{(i,j,r)}\exp\!\bigl(-(y/\omega_{ij,r})^{-\theta}\bigr)
  \;=\;\exp\!\bigl(-y^{-\theta}S\bigr),
  \qquad
  S\;\equiv\;\sum_{(i,j,r)}\omega_{ij,r}^{\theta}.
\]
Using \eqref{eq:tau_def} we have $S=\sum_{(i,j)}u_i^{\theta}w_j^{\theta}\tijk^{-\theta}
=\Wbar^{\theta}$, so $M$ is Fr\'echet with scale $\Wbar$ and
$\mathbb E[M]=\Gamma(1-1/\theta)\Wbar$, giving
\eqref{eq:welfare_main}.

The ex-post indifference property follows from $\mathbb E[V\mid(i,j)
\text{ chosen}]=\mathbb E[M]$: after conditioning on the chosen pair,
the distribution of $V$ is the same Fr\'echet upper tail regardless of
$(i,j)$.

\subsection{Link-traffic gravity}\label{app:linkgravity}

For a worker choosing the pair $(i,j)$, the expected number of times
edge $(k,l)$ is traversed along her chosen route equals the total
$\theta$-weighted walk-mass through $(k,l)$ divided by the total
$\theta$-weighted walk-mass from $i$ to $j$:
\[
  \pi^{kl}_{ij}
  \;=\;
  \frac{\tau_{ik}^{-\theta}\,\tkl^{-\theta}\,\tau_{lj}^{-\theta}}
       {\tijk^{-\theta}}.
\]
Aggregating with the gravity flows \eqref{eq:gravity_MA},
\[
  \xi_{kl}
  \;=\;\sum_{i,j}L_{ij}\,\pi^{kl}_{ij}
  \;=\;\tkl^{-\theta}\sum_{i}\tau_{ik}^{-\theta}L_i^R\Pi_i^{\theta}\cdot
                     \sum_{j}\tau_{lj}^{-\theta}L_j^FP_j^{\theta}
  \;=\;\tkl^{-\theta}\,P_k^{-\theta}\,\Pi_l^{-\theta},
\]
where the last equality invokes \eqref{eq:MA_def}.  This proves
Proposition~\ref{prop:linkgravity}.

\subsection{Equilibrium system in shares}\label{app:eqshares}

Substituting \eqref{eq:fundamentals} and \eqref{eq:gravity} into
$L_i^R=\sum_jL_{ij}$ and dividing through by $\Lbar$ yields
\eqref{eq:eq_R} with
$\chi=\Lbar^{\theta(\alpha+\beta)}/\Wbar^{\theta}$; the argument for
\eqref{eq:eq_F} is symmetric.  Existence of a strictly positive
solution follows from the Brouwer-based fixed-point construction in
\citet{AllenArkolakis2022}; uniqueness under \eqref{eq:unique}
follows from the bounding argument in their Proposition~1(b).
Numerical verification of \eqref{eq:unique} over the parameter grid
in \S\ref{ssec:data}, together with spectral-radius checks on the
Jacobian of the iteration \eqref{eq:invert}, confirms convergence to
a unique strictly positive equilibrium for every parameter
constellation considered in the paper.

\subsection{Existence and uniqueness, proof sketch}\label{app:exist}

Let $F:(l^R,l^F)\mapsto(\hat l^R,\hat l^F)$ denote the right-hand side
of \eqref{eq:eq_R}--\eqref{eq:eq_F} normalised by $\chi$.  $F$ is a
continuous positive operator on the simplex $\{(l^R,l^F):\sum l^R=
\sum l^F=1\}$.  Brouwer's theorem delivers existence of a fixed
point, and strict positivity follows from the strict positivity of
the primitives.  For uniqueness, write $F$ in log-linear form and
bound the spectral norm of the Jacobian under \eqref{eq:unique}; the
bound specialises to Proposition~1(b) of \citet{AllenArkolakis2022}
when restricted to the urban variant.  Global contractivity is not
established in closed form but is verified numerically at every
parameter combination used in the paper.

\subsection{Hat algebra}\label{app:hat}

Taking hats in \eqref{eq:eq_R}--\eqref{eq:eq_F} and using
$\hat u_i=(\hat l_i^R)^{\beta}$, $\hat w_j=(\hat l_j^F)^{\alpha}$
delivers \eqref{eq:hat_system}.  The coupling with the congestion
technology is
\[
  \hat\xi_{kl}\;=\;\hat\tkl^{-\theta}\hat P_k^{-\theta}\hat\Pi_l^{-\theta},
  \qquad
  \hat\tkl^{1+\theta\lambda}
  \;=\;\hat{\tbar}_{kl}\,\hat n_{kl}^{-\lambda}\,(\hat P_k\hat\Pi_l)^{-\theta\lambda}.
\]
We solve the full coupled hat system by a nested fixed-point
procedure: the outer loop iterates on $\hat\tkl$, and the inner loop
solves \eqref{eq:hat_system} for $(\hat l_i^R,\hat l_j^F)$ given
$\{\hat\tkl\}$.  Welfare follows from \eqref{eq:welfare_hat}.

\subsection{Data}\label{app:data}

The sample is restricted to the largest weakly connected component
of $(\N,\E)$.  The raw Baidu commuting matrix contains 442{,}796
distinct origin--destination pairs and $7.96$ million AM-peak commute
observations; after spatial and reachability filters the estimation
sample is $N=228$ grid cells and 704 directed links.  Descriptive
statistics on $(\tbar_{kl},\tkl^{\obs},n_{kl})$ and diagnostics on
the coverage of the matched centreline dataset are available upon
request.

\subsection{Identification of $\lambda$}\label{app:iv}

The first stage of the IV procedure projects $\ln(\xi_{kl}^{\obs}/n_{kl})$
on the instrument (count of signalised intersections along the link)
conditional on origin and destination fixed effects and on bilateral
route-quality controls.  Identification requires that, conditional on
fixed effects, the instrument affects the equilibrium traffic $\xi_{kl}$
through the demand side of the link-traffic gravity equation
\eqref{eq:xi_gravity} but is uncorrelated with the free-flow component
$\tbar_{kl}$ of the travel-time cost.  First-stage F-statistics
exceed 30, weak-instrument concerns are rejected at conventional
thresholds, and the baseline 2SLS estimate is robust to alternative
instrument constructions; full tables and falsification exercises are
in the empirical supplement.

%====================================================================
%  Bibliography
%====================================================================
\begin{thebibliography}{10}
\bibitem[Ahlfeldt et al.(2015)]{Ahlfeldt2015}
Ahlfeldt, G., Redding, S. J., Sturm, D. M., Wolf, N. (2015).
The economics of density: Evidence from the Berlin Wall.
\emph{Econometrica}, 83(6), 2127--2189.

\bibitem[Allen and Arkolakis(2014)]{AllenArkolakis2014}
Allen, T., Arkolakis, C. (2014). Trade and the topography of the
spatial economy. \emph{Quarterly Journal of Economics}, 129(3),
1085--1140.

\bibitem[Allen and Arkolakis(2022)]{AllenArkolakis2022}
Allen, T., Arkolakis, C. (2022). The welfare effects of transportation
infrastructure improvements. \emph{Review of Economic Studies},
89(6), 2911--2957.

\bibitem[Combes and Gobillon(2015)]{CombesGobillon2015}
Combes, P.-P., Gobillon, L. (2015). The empirics of agglomeration
economies. In \emph{Handbook of Regional and Urban Economics},
Vol.~5A, 247--348. Elsevier.

\bibitem[Duranton and Turner(2011)]{DurantonTurner2011}
Duranton, G., Turner, M. A. (2011). The fundamental law of road
congestion: Evidence from US cities. \emph{American Economic Review},
101(6), 2616--2652.

\bibitem[Heblich, Redding and Sturm(2020)]{HeblichReddingSturm2020}
Heblich, S., Redding, S. J., Sturm, D. M. (2020). The making of the
modern metropolis: Evidence from London. \emph{Quarterly Journal of
Economics}, 135(4), 2059--2133.

\bibitem[Monte, Redding and Rossi-Hansberg(2018)]{MonteReddingRossiHansberg2018}
Monte, F., Redding, S. J., Rossi-Hansberg, E. (2018). Commuting,
migration, and local employment elasticities. \emph{American Economic
Review}, 108(12), 3855--3890.

\bibitem[Rosenthal and Strange(2004)]{RosenthalStrange2004}
Rosenthal, S. S., Strange, W. C. (2004). Evidence on the nature and
sources of agglomeration economies. In \emph{Handbook of Regional and
Urban Economics}, Vol.~4, 2119--2171. Elsevier.
\end{thebibliography}

\end{document}
